
\documentclass{article}
\usepackage{colortbl}
\usepackage{makecell}
\usepackage{multirow}
\usepackage{supertabular}

\begin{document}

\newcounter{utterance}

\centering \large Interaction Transcript for game `clean\_up', experiment `2\_hard\_5obj\_lv', episode 3 with glm{-}5.2.
\vspace{24pt}

{ \footnotesize  \setcounter{utterance}{1}
\setlength{\tabcolsep}{0pt}
\begin{supertabular}{c@{$\;$}|p{.15\linewidth}@{}p{.15\linewidth}p{.15\linewidth}p{.15\linewidth}p{.15\linewidth}p{.15\linewidth}}
    \# & $\;$A & \multicolumn{4}{c}{Game Master} & $\;\:$B\\
    \hline

    \theutterance \stepcounter{utterance}  
    & & \multicolumn{4}{p{0.6\linewidth}}{
        \cellcolor[rgb]{0.9,0.9,0.9}{
            \makecell[{{p{\linewidth}}}]{
                \texttt{\tiny{[P1$\langle$GM]}}
                \texttt{Es esmu jūsu spēles vadītājs, un jūs spēlējat sadarbības spēli ar šādu režģi kā spēles laukumu:} \\
\\ 
\texttt{```} \\
\texttt{ 1234567} \\
\texttt{╔═╤═╤═══╗} \\
\texttt{║C│◌│◌TI║ 1} \\
\texttt{║◌├─┤W◌◌║ 2} \\
\texttt{╟─┤◌├─┬─╢ 3} \\
\texttt{║◌│◌│◌│◌║ 4} \\
\texttt{╟─┼─┘◌│◌║ 5} \\
\texttt{║◌│◌H┌┴─╢ 6} \\
\texttt{║◌│◌◌│◌◌║ 7} \\
\texttt{╚═╧══╧══╝} \\
\\ 
\texttt{```} \\
\\ 
\texttt{* Augšējā malā ir attēlotas x koordinātas, kas pieaug pa labi, un labajā malā y koordinātas, kas pieaug uz leju.} \\
\texttt{* Uz jūsu režģa ir nejauši izvietoti šādi objekti: 'W', 'I', 'T', 'C', 'H'.} \\
\\ 
\\ 
\texttt{Otrs spēlētājs redz spēles laukuma variāciju, kurā objekti ir izvietoti citās nejaušās vietās. Jūs nevarat redzēt otra spēlētāja laukumu, un viņš nevar redzēt jūsējo.} \\
\\ 
\\ 
\texttt{**Mērķis:**} \\
\\ 
\texttt{Abiem spēlētājiem ir jāpārvieto objekti uz sava fona tā, lai vienādi objekti nonāktu vienādās koordinātēs. Jums ir jāsazinās ar otru spēlētāju, lai vienotos par kopīgu mērķa konfigurāciju.} \\
\\ 
\\ 
\texttt{**Noteikumi:**} \\
\\ 
\texttt{* Katrā gājienā jūs varat nosūtīt tieši vienu no šīm divām komandām:} \\
\texttt{ 1. `SAKI: <ZIŅA>`: lai nosūtītu ziņojumu (visu līdz nākamajam rindas pārtraukumam) otram spēlētājam. Es to pārsūtīšu jūsu partnerim.} \\
\texttt{2. `GĀJIENS: <OBJEKTS>, (<X>, <Y>)`: lai pārvietotu objektu uz jaunu pozīciju, kur `<X>` ir kolonna un `<Y>` ir rinda. Es jūs informēšu, vai jūsu gājiens bija veiksmīgs vai nē.} \\
\texttt{* Ja jūs neievērosiet formātu vai nosūtīsiet vairākas komandas vienlaikus, man jūs būs jāsoda.} \\
\texttt{* Ja abi spēlētāji kopā uzkrāj vairāk nekā 12 sodus, jūs abi zaudējat spēli.} \\
\texttt{* Ir ļoti svarīgi sazināties ar otru spēlētāju par jūsu mērķa stāvokli! *Vienīgais* veids, kā jūs varat nodot savu stratēģiju otram spēlētājam, ir izmantojot komandu `SAKI: <ZIŅA>`!} \\
\\ 
\\ 
\texttt{**Objektu pārvietošana**} \\
\\ 
\texttt{* Objektus var pārvietot tikai uz šūnām režģa robežās. Mērķa šūnai jābūt tukšai, t.i., tajā jābūt tikai simbolam '◌'.} \\
\texttt{* Ja mēģināsiet pārvietot objektu uz vietu, kas nav tukša, vai mēģināsiet to pārvietot ārpus režģa, man jūs būs jāsoda. Jūs saņemat vēl vienu mēģinājumu.} \\
\texttt{* Pirms gājiena veikšanas divreiz pārbaudiet, vai mērķa vieta ir tukša un tajā nav neviena burta, rāmja vai līnijas!} \\
\\ 
\\ 
\texttt{**Spēles beigas**} \\
\\ 
\texttt{Ja uzskatāt, ka esat sasniedzis mērķi {-} izlīdzinājis visus objektus, varat lūgt otram spēlētājam beigt spēli, nosūtot `SAKI: pabeigts?`. Ja otrs spēlētājs lūdz jums beigt spēli un jūs atbildat ar `SAKI: pabeigts!`, spēle beigsies.} \\
\\ 
\texttt{Abi spēlētāji uzvar, ja spēle beidzas 20 raundos, kur viens raunds ir definēts kā divi spēlētāji, kas katrs nosūta derīgu komandu.} \\
\\ 
\texttt{**Punktu skaitīšana:**} \\
\\ 
\texttt{Jo tuvāk vienādi objekti atrodas abos spēles laukumos, jo vairāk punktu jūs saņemat. Sodi samazina jūsu punktus. Vai varat pārspēt rekordu?} \\
\\ 
\\ 
\texttt{Lūdzu, nosūtiet ziņu otram spēlētājam, lai sāktu spēli!} \\
            }
        }
    }
    & \\ \\

    \theutterance \stepcounter{utterance}  
    & \multicolumn{4}{p{0.6\linewidth}}{
        \cellcolor[rgb]{0.9,0.9,0.9}{
            \makecell[{{p{\linewidth}}}]{
                \texttt{\tiny{[P1$\rangle$GM]}}
                \texttt{Sveiks! Sāksim spēli. Lai sasniegtu mērķi, mums jāvienojas par kopīgu mērķa konfigurāciju. Šeit ir mans priekšlikums:} \\
\\ 
\texttt{Mans piedāvājums:} \\
\texttt{C: (1, 4)} \\
\texttt{W: (4, 6)} \\
\texttt{I: (4, 7)} \\
\texttt{T: (5, 1)} \\
\texttt{H: (3, 4)} \\
\\ 
\texttt{Vai šī konfigurācija tev ir piemērota? Ja ir, sāksim pārvietot objektus. Ja nē, lūdzu, piedāvā savu variantu!} \\
\\ 
\texttt{SAKI: Sveiki! Mans priekšlikums mērķa koordinātēm: C(1,4), W(4,6), I(4,7), T(5,1), H(3,4). Vai piekrīti?} \\
            }
        }
    }
    & & \\ \\

    \theutterance \stepcounter{utterance}  
    & & & \multicolumn{2}{p{0.3\linewidth}}{
        \cellcolor[rgb]{0.9,0.9,0.9}{
            \makecell[{{p{\linewidth}}}]{
                \texttt{\tiny{[GM$|$GM]}}
                \texttt{Sveiks! Sāksim spēli. Lai sasniegtu mērķi, mums jāvienojas par kopīgu mērķa konfigurāciju. Šeit ir mans priekšlikums:} \\
\\ 
\texttt{Mans piedāvājums:} \\
\texttt{C: (1, 4)} \\
\texttt{W: (4, 6)} \\
\texttt{I: (4, 7)} \\
\texttt{T: (5, 1)} \\
\texttt{H: (3, 4)} \\
\\ 
\texttt{Vai šī konfigurācija tev ir piemērota? Ja ir, sāksim pārvietot objektus. Ja nē, lūdzu, piedāvā savu variantu!} \\
\\ 
\texttt{SAKI: Sveiki! Mans priekšlikums mērķa koordinātēm: C(1,4), W(4,6), I(4,7), T(5,1), H(3,4). Vai piekrīti?} \\
            }
        }
    }
    & & \\ \\

    \theutterance \stepcounter{utterance}  
    & & & \multicolumn{2}{p{0.3\linewidth}}{
        \cellcolor[rgb]{0.9,0.9,0.9}{
            \makecell[{{p{\linewidth}}}]{
                \texttt{\tiny{[GM$|$GM]}}
                \texttt{Invalid format: head is not empty: 'Sveiks! Sāksim spēli. Lai sasniegtu mērķi, mums jāvienojas par kopīgu mērķa konfigurāciju. Šeit ir mans priekšlikums:} \\
\\ 
\texttt{Mans piedāvājums:} \\
\texttt{C: (1, 4)} \\
\texttt{W: (4, 6)} \\
\texttt{I: (4, 7)} \\
\texttt{T: (5, 1)} \\
\texttt{H: (3, 4)} \\
\\ 
\texttt{Vai šī konfigurācija tev ir piemērota? Ja ir, sāksim pārvietot objektus. Ja nē, lūdzu, piedāvā savu variantu!} \\
\\ 
\texttt{'} \\
            }
        }
    }
    & & \\ \\

    \theutterance \stepcounter{utterance}  
    & & \multicolumn{4}{p{0.6\linewidth}}{
        \cellcolor[rgb]{0.9,0.9,0.9}{
            \makecell[{{p{\linewidth}}}]{
                \texttt{\tiny{[P1$\langle$GM]}}
                \texttt{Sods: Jūsu ziņojumā nedrīkst būt nekas pirms komandas!} \\
\\ 
\texttt{Pārliecinieties, ka jūsu atbildē ir tikai `SAKI: <ZIŅA>` vai `GĀJIENS: <OBJEKTS>, (<X>, <Y>)`, un nekas cits!} \\
\\ 
\texttt{Jūs kopā esat uzkrājuši 1 no 12 sodiem.} \\
\texttt{Lūdzu, mēģiniet vēlreiz!} \\
            }
        }
    }
    & \\ \\

    \theutterance \stepcounter{utterance}  
    & \multicolumn{4}{p{0.6\linewidth}}{
        \cellcolor[rgb]{0.9,0.9,0.9}{
            \makecell[{{p{\linewidth}}}]{
                \texttt{\tiny{[P1$\rangle$GM]}}
                \texttt{SAKI: Sveiki! Mans priekšlikums mērķa koordinātēm: C(1,4), W(4,6), I(4,7), T(5,1), H(3,4). Vai piekrīti?} \\
            }
        }
    }
    & & \\ \\

    \theutterance \stepcounter{utterance}  
    & & & \multicolumn{2}{p{0.3\linewidth}}{
        \cellcolor[rgb]{0.9,0.9,0.9}{
            \makecell[{{p{\linewidth}}}]{
                \texttt{\tiny{[GM$|$GM]}}
                \texttt{SAKI: Sveiki! Mans priekšlikums mērķa koordinātēm: C(1,4), W(4,6), I(4,7), T(5,1), H(3,4). Vai piekrīti?} \\
            }
        }
    }
    & & \\ \\

    \theutterance \stepcounter{utterance}  
    & & \multicolumn{4}{p{0.6\linewidth}}{
        \cellcolor[rgb]{0.9,0.9,0.9}{
            \makecell[{{p{\linewidth}}}]{
                \texttt{\tiny{[GM$\rangle$P2]}}
                \texttt{Es esmu jūsu spēles vadītājs, un jūs spēlējat sadarbības spēli ar šādu režģi kā spēles laukumu:} \\
\\ 
\texttt{```} \\
\texttt{ 1234567} \\
\texttt{╔═╤═╤═══╗} \\
\texttt{║H│◌│◌◌◌║ 1} \\
\texttt{║◌├─┤◌◌◌║ 2} \\
\texttt{╟─┤◌├─┬─╢ 3} \\
\texttt{║◌│◌│C│◌║ 4} \\
\texttt{╟─┼─┘◌│◌║ 5} \\
\texttt{║◌│◌◌┌┴─╢ 6} \\
\texttt{║T│WI│◌◌║ 7} \\
\texttt{╚═╧══╧══╝} \\
\\ 
\texttt{```} \\
\\ 
\texttt{* Augšējā malā ir attēlotas x koordinātas, kas pieaug pa labi, un labajā malā y koordinātas, kas pieaug uz leju.} \\
\texttt{* Uz jūsu režģa ir nejauši izvietoti šādi objekti: 'W', 'I', 'T', 'C', 'H'.} \\
\\ 
\\ 
\texttt{Otrs spēlētājs redz spēles laukuma variāciju, kurā objekti ir izvietoti citās nejaušās vietās. Jūs nevarat redzēt otra spēlētāja laukumu, un viņš nevar redzēt jūsējo.} \\
\\ 
\\ 
\texttt{**Mērķis:**} \\
\\ 
\texttt{Abiem spēlētājiem ir jāpārvieto objekti uz sava fona tā, lai vienādi objekti nonāktu vienādās koordinātēs. Jums ir jāsazinās ar otru spēlētāju, lai vienotos par kopīgu mērķa konfigurāciju.} \\
\\ 
\\ 
\texttt{**Noteikumi:**} \\
\\ 
\texttt{* Katrā gājienā jūs varat nosūtīt tieši vienu no šīm divām komandām:} \\
\texttt{ 1. `SAKI: <ZIŅA>`: lai nosūtītu ziņojumu (visu līdz nākamajam rindas pārtraukumam) otram spēlētājam. Es to pārsūtīšu jūsu partnerim.} \\
\texttt{2. `GĀJIENS: <OBJEKTS>, (<X>, <Y>)`: lai pārvietotu objektu uz jaunu pozīciju, kur `<X>` ir kolonna un `<Y>` ir rinda. Es jūs informēšu, vai jūsu gājiens bija veiksmīgs vai nē.} \\
\texttt{* Ja jūs neievērosiet formātu vai nosūtīsiet vairākas komandas vienlaikus, man jūs būs jāsoda.} \\
\texttt{* Ja abi spēlētāji kopā uzkrāj vairāk nekā 12 sodus, jūs abi zaudējat spēli.} \\
\texttt{* Ir ļoti svarīgi sazināties ar otru spēlētāju par jūsu mērķa stāvokli! *Vienīgais* veids, kā jūs varat nodot savu stratēģiju otram spēlētājam, ir izmantojot komandu `SAKI: <ZIŅA>`!} \\
\\ 
\\ 
\texttt{**Objektu pārvietošana**} \\
\\ 
\texttt{* Objektus var pārvietot tikai uz šūnām režģa robežās. Mērķa šūnai jābūt tukšai, t.i., tajā jābūt tikai simbolam '◌'.} \\
\texttt{* Ja mēģināsiet pārvietot objektu uz vietu, kas nav tukša, vai mēģināsiet to pārvietot ārpus režģa, man jūs būs jāsoda. Jūs saņemat vēl vienu mēģinājumu.} \\
\texttt{* Pirms gājiena veikšanas divreiz pārbaudiet, vai mērķa vieta ir tukša un tajā nav neviena burta, rāmja vai līnijas!} \\
\\ 
\\ 
\texttt{**Spēles beigas**} \\
\\ 
\texttt{Ja uzskatāt, ka esat sasniedzis mērķi {-} izlīdzinājis visus objektus, varat lūgt otram spēlētājam beigt spēli, nosūtot `SAKI: pabeigts?`. Ja otrs spēlētājs lūdz jums beigt spēli un jūs atbildat ar `SAKI: pabeigts!`, spēle beigsies.} \\
\\ 
\texttt{Abi spēlētāji uzvar, ja spēle beidzas 20 raundos, kur viens raunds ir definēts kā divi spēlētāji, kas katrs nosūta derīgu komandu.} \\
\\ 
\texttt{**Punktu skaitīšana:**} \\
\\ 
\texttt{Jo tuvāk vienādi objekti atrodas abos spēles laukumos, jo vairāk punktu jūs saņemat. Sodi samazina jūsu punktus. Vai varat pārspēt rekordu?} \\
\\ 
\texttt{Otrs spēlētājs sāka spēli, nosūtot šo ziņojumu:} \\
\\ 
\texttt{"Sveiki! Mans priekšlikums mērķa koordinātēm: C(1,4), W(4,6), I(4,7), T(5,1), H(3,4). Vai piekrīti?"} \\
\\ 
\texttt{Kāda ir jūsu pirmā komanda?} \\
            }
        }
    }
    & \\ \\

    \theutterance \stepcounter{utterance}  
    & & & \multicolumn{4}{p{0.6\linewidth}}{
        \cellcolor[rgb]{0.9,0.9,0.9}{
            \makecell[{{p{\linewidth}}}]{
                \texttt{\tiny{[GM$\langle$P2]}}
                \texttt{SAKI: Sveiks! Piekrītu tavam priekšlikumam: C(1,4), W(4,6), I(4,7), T(5,1), H(3,4). Šīs pozīcijas man der, sāksim pārvietošanu.} \\
            }
        }
    }
     \\ \\

    \theutterance \stepcounter{utterance}  
    & & & \multicolumn{2}{p{0.3\linewidth}}{
        \cellcolor[rgb]{0.9,0.9,0.9}{
            \makecell[{{p{\linewidth}}}]{
                \texttt{\tiny{[GM$|$GM]}}
                \texttt{SAKI: Sveiks! Piekrītu tavam priekšlikumam: C(1,4), W(4,6), I(4,7), T(5,1), H(3,4). Šīs pozīcijas man der, sāksim pārvietošanu.} \\
            }
        }
    }
    & & \\ \\

    \theutterance \stepcounter{utterance}  
    & & \multicolumn{4}{p{0.6\linewidth}}{
        \cellcolor[rgb]{0.9,0.9,0.9}{
            \makecell[{{p{\linewidth}}}]{
                \texttt{\tiny{[P1$\langle$GM]}}
                \texttt{Jūsu ziņojums ir nodots otram spēlētājam.} \\
\\ 
\\ 
\texttt{Pašlaik jūs spēlējat kārtu 1 no maksimāli 20.} \\
\texttt{Jūs kopā esat uzkrājuši 1 no 12 sodiem.} \\
\texttt{Otrs spēlētājs nosūtīja šo ziņojumu:} \\
\\ 
\\ 
\texttt{"Sveiks! Piekrītu tavam priekšlikumam: C(1,4), W(4,6), I(4,7), T(5,1), H(3,4). Šīs pozīcijas man der, sāksim pārvietošanu."} \\
\\ 
\\ 
\texttt{Kāda ir jūsu nākamā komanda?} \\
            }
        }
    }
    & \\ \\

    \theutterance \stepcounter{utterance}  
    & \multicolumn{4}{p{0.6\linewidth}}{
        \cellcolor[rgb]{0.9,0.9,0.9}{
            \makecell[{{p{\linewidth}}}]{
                \texttt{\tiny{[P1$\rangle$GM]}}
                \texttt{GĀJIENS: T, (4, 1)} \\
            }
        }
    }
    & & \\ \\

    \theutterance \stepcounter{utterance}  
    & & & \multicolumn{2}{p{0.3\linewidth}}{
        \cellcolor[rgb]{0.9,0.9,0.9}{
            \makecell[{{p{\linewidth}}}]{
                \texttt{\tiny{[GM$|$GM]}}
                \texttt{GĀJIENS: T, (4, 1)} \\
            }
        }
    }
    & & \\ \\

    \theutterance \stepcounter{utterance}  
    & & & \multicolumn{2}{p{0.3\linewidth}}{
        \cellcolor[rgb]{0.9,0.9,0.9}{
            \makecell[{{p{\linewidth}}}]{
                \texttt{\tiny{[GM$|$GM]}}
                \texttt{Sods: (4,1) nav tukšs, bet satur '│'.} \\
\\ 
\texttt{Jūs kopā esat uzkrājuši 2 no 12 sodiem.} \\
\texttt{Lūdzu, mēģiniet vēlreiz!} \\
            }
        }
    }
    & & \\ \\

    \theutterance \stepcounter{utterance}  
    & & \multicolumn{4}{p{0.6\linewidth}}{
        \cellcolor[rgb]{0.9,0.9,0.9}{
            \makecell[{{p{\linewidth}}}]{
                \texttt{\tiny{[P1$\langle$GM]}}
                \texttt{Sods: (4,1) nav tukšs, bet satur '│'.} \\
\\ 
\texttt{Jūs kopā esat uzkrājuši 2 no 12 sodiem.} \\
\texttt{Lūdzu, mēģiniet vēlreiz!} \\
            }
        }
    }
    & \\ \\

    \theutterance \stepcounter{utterance}  
    & \multicolumn{4}{p{0.6\linewidth}}{
        \cellcolor[rgb]{0.9,0.9,0.9}{
            \makecell[{{p{\linewidth}}}]{
                \texttt{\tiny{[P1$\rangle$GM]}}
                \texttt{Looking at the grid more carefully, (5,1) contains 'T' already, so the target for T should be a nearby empty cell. Let me move T to (3,1) which appears to be empty (◌).} \\
\\ 
\texttt{GĀJIENS: T, (3, 1)} \\
            }
        }
    }
    & & \\ \\

    \theutterance \stepcounter{utterance}  
    & & & \multicolumn{2}{p{0.3\linewidth}}{
        \cellcolor[rgb]{0.9,0.9,0.9}{
            \makecell[{{p{\linewidth}}}]{
                \texttt{\tiny{[GM$|$GM]}}
                \texttt{Looking at the grid more carefully, (5,1) contains 'T' already, so the target for T should be a nearby empty cell. Let me move T to (3,1) which appears to be empty (◌).} \\
\\ 
\texttt{GĀJIENS: T, (3, 1)} \\
            }
        }
    }
    & & \\ \\

    \theutterance \stepcounter{utterance}  
    & & & \multicolumn{2}{p{0.3\linewidth}}{
        \cellcolor[rgb]{0.9,0.9,0.9}{
            \makecell[{{p{\linewidth}}}]{
                \texttt{\tiny{[GM$|$GM]}}
                \texttt{Invalid format: head is not empty: 'Looking at the grid more carefully, (5,1) contains 'T' already, so the target for T should be a nearby empty cell. Let me move T to (3,1) which appears to be empty (◌).} \\
\\ 
\texttt{'} \\
            }
        }
    }
    & & \\ \\

    \theutterance \stepcounter{utterance}  
    & & \multicolumn{4}{p{0.6\linewidth}}{
        \cellcolor[rgb]{0.9,0.9,0.9}{
            \makecell[{{p{\linewidth}}}]{
                \texttt{\tiny{[P1$\langle$GM]}}
                \texttt{Sods: Jūsu ziņojumā nedrīkst būt nekas pirms komandas!} \\
\\ 
\texttt{Pārliecinieties, ka jūsu atbildē ir tikai `SAKI: <ZIŅA>` vai `GĀJIENS: <OBJEKTS>, (<X>, <Y>)`, un nekas cits!} \\
\\ 
\texttt{Jūs kopā esat uzkrājuši 3 no 12 sodiem.} \\
\texttt{Lūdzu, mēģiniet vēlreiz!} \\
            }
        }
    }
    & \\ \\

    \theutterance \stepcounter{utterance}  
    & \multicolumn{4}{p{0.6\linewidth}}{
        \cellcolor[rgb]{0.9,0.9,0.9}{
            \makecell[{{p{\linewidth}}}]{
                \texttt{\tiny{[P1$\rangle$GM]}}
                \texttt{GĀJIENS: T, (3, 1)} \\
            }
        }
    }
    & & \\ \\

    \theutterance \stepcounter{utterance}  
    & & & \multicolumn{2}{p{0.3\linewidth}}{
        \cellcolor[rgb]{0.9,0.9,0.9}{
            \makecell[{{p{\linewidth}}}]{
                \texttt{\tiny{[GM$|$GM]}}
                \texttt{GĀJIENS: T, (3, 1)} \\
            }
        }
    }
    & & \\ \\

    \theutterance \stepcounter{utterance}  
    & & & \multicolumn{2}{p{0.3\linewidth}}{
        \cellcolor[rgb]{0.9,0.9,0.9}{
            \makecell[{{p{\linewidth}}}]{
                \texttt{\tiny{[GM$|$GM]}}
                \texttt{T moved to (3, 1)} \\
            }
        }
    }
    & & \\ \\

    \theutterance \stepcounter{utterance}  
    & & \multicolumn{4}{p{0.6\linewidth}}{
        \cellcolor[rgb]{0.9,0.9,0.9}{
            \makecell[{{p{\linewidth}}}]{
                \texttt{\tiny{[GM$\rangle$P2]}}
                \texttt{Jūsu ziņojums ir nodots otram spēlētājam.} \\
\\ 
\\ 
\texttt{Pašlaik jūs spēlējat kārtu 2 no maksimāli 20.} \\
\texttt{Jūs kopā esat uzkrājuši 3 no 12 sodiem.} \\
\texttt{Otrs spēlētājs pārvietoja objektu savā režģī.} \\
\\ 
\\ 
\texttt{Kāda ir jūsu nākamā komanda?} \\
            }
        }
    }
    & \\ \\

    \theutterance \stepcounter{utterance}  
    & & & \multicolumn{4}{p{0.6\linewidth}}{
        \cellcolor[rgb]{0.9,0.9,0.9}{
            \makecell[{{p{\linewidth}}}]{
                \texttt{\tiny{[GM$\langle$P2]}}
                \texttt{GĀJIENS: T, (5,1)} \\
            }
        }
    }
     \\ \\

    \theutterance \stepcounter{utterance}  
    & & & \multicolumn{2}{p{0.3\linewidth}}{
        \cellcolor[rgb]{0.9,0.9,0.9}{
            \makecell[{{p{\linewidth}}}]{
                \texttt{\tiny{[GM$|$GM]}}
                \texttt{GĀJIENS: T, (5,1)} \\
            }
        }
    }
    & & \\ \\

    \theutterance \stepcounter{utterance}  
    & & & \multicolumn{2}{p{0.3\linewidth}}{
        \cellcolor[rgb]{0.9,0.9,0.9}{
            \makecell[{{p{\linewidth}}}]{
                \texttt{\tiny{[GM$|$GM]}}
                \texttt{T moved to (5, 1)} \\
            }
        }
    }
    & & \\ \\

    \theutterance \stepcounter{utterance}  
    & & \multicolumn{4}{p{0.6\linewidth}}{
        \cellcolor[rgb]{0.9,0.9,0.9}{
            \makecell[{{p{\linewidth}}}]{
                \texttt{\tiny{[P1$\langle$GM]}}
                \texttt{Veiksmīgi pārvietots 'T' uz (3,1). Jūsu atjauninātais režģis izskatās šādi:} \\
\texttt{```} \\
\texttt{ 1234567} \\
\texttt{╔═╤═╤═══╗} \\
\texttt{║C│T│◌◌I║ 1} \\
\texttt{║◌├─┤W◌◌║ 2} \\
\texttt{╟─┤◌├─┬─╢ 3} \\
\texttt{║◌│◌│◌│◌║ 4} \\
\texttt{╟─┼─┘◌│◌║ 5} \\
\texttt{║◌│◌H┌┴─╢ 6} \\
\texttt{║◌│◌◌│◌◌║ 7} \\
\texttt{╚═╧══╧══╝} \\
\texttt{```} \\
\texttt{Pašlaik jūs spēlējat kārtu 2 no maksimāli 20.} \\
\texttt{Jūs kopā esat uzkrājuši 3 no 12 sodiem.} \\
\texttt{Otrs spēlētājs pārvietoja objektu savā režģī.} \\
\\ 
\\ 
\texttt{Kāda ir jūsu nākamā komanda?} \\
            }
        }
    }
    & \\ \\

    \theutterance \stepcounter{utterance}  
    & \multicolumn{4}{p{0.6\linewidth}}{
        \cellcolor[rgb]{0.9,0.9,0.9}{
            \makecell[{{p{\linewidth}}}]{
                \texttt{\tiny{[P1$\rangle$GM]}}
                \texttt{GĀJIENS: C, (1, 4)} \\
            }
        }
    }
    & & \\ \\

    \theutterance \stepcounter{utterance}  
    & & & \multicolumn{2}{p{0.3\linewidth}}{
        \cellcolor[rgb]{0.9,0.9,0.9}{
            \makecell[{{p{\linewidth}}}]{
                \texttt{\tiny{[GM$|$GM]}}
                \texttt{GĀJIENS: C, (1, 4)} \\
            }
        }
    }
    & & \\ \\

    \theutterance \stepcounter{utterance}  
    & & & \multicolumn{2}{p{0.3\linewidth}}{
        \cellcolor[rgb]{0.9,0.9,0.9}{
            \makecell[{{p{\linewidth}}}]{
                \texttt{\tiny{[GM$|$GM]}}
                \texttt{C moved to (1, 4)} \\
            }
        }
    }
    & & \\ \\

    \theutterance \stepcounter{utterance}  
    & & \multicolumn{4}{p{0.6\linewidth}}{
        \cellcolor[rgb]{0.9,0.9,0.9}{
            \makecell[{{p{\linewidth}}}]{
                \texttt{\tiny{[GM$\rangle$P2]}}
                \texttt{Veiksmīgi pārvietots 'T' uz (5,1). Jūsu atjauninātais režģis izskatās šādi:} \\
\texttt{```} \\
\texttt{ 1234567} \\
\texttt{╔═╤═╤═══╗} \\
\texttt{║H│◌│T◌◌║ 1} \\
\texttt{║◌├─┤◌◌◌║ 2} \\
\texttt{╟─┤◌├─┬─╢ 3} \\
\texttt{║◌│◌│C│◌║ 4} \\
\texttt{╟─┼─┘◌│◌║ 5} \\
\texttt{║◌│◌◌┌┴─╢ 6} \\
\texttt{║◌│WI│◌◌║ 7} \\
\texttt{╚═╧══╧══╝} \\
\texttt{```} \\
\texttt{Pašlaik jūs spēlējat kārtu 3 no maksimāli 20.} \\
\texttt{Jūs kopā esat uzkrājuši 3 no 12 sodiem.} \\
\texttt{Otrs spēlētājs pārvietoja objektu savā režģī.} \\
\\ 
\\ 
\texttt{Kāda ir jūsu nākamā komanda?} \\
            }
        }
    }
    & \\ \\

    \theutterance \stepcounter{utterance}  
    & & & \multicolumn{4}{p{0.6\linewidth}}{
        \cellcolor[rgb]{0.9,0.9,0.9}{
            \makecell[{{p{\linewidth}}}]{
                \texttt{\tiny{[GM$\langle$P2]}}
                \texttt{GĀJIENS: H, (3,4)} \\
            }
        }
    }
     \\ \\

    \theutterance \stepcounter{utterance}  
    & & & \multicolumn{2}{p{0.3\linewidth}}{
        \cellcolor[rgb]{0.9,0.9,0.9}{
            \makecell[{{p{\linewidth}}}]{
                \texttt{\tiny{[GM$|$GM]}}
                \texttt{GĀJIENS: H, (3,4)} \\
            }
        }
    }
    & & \\ \\

    \theutterance \stepcounter{utterance}  
    & & & \multicolumn{2}{p{0.3\linewidth}}{
        \cellcolor[rgb]{0.9,0.9,0.9}{
            \makecell[{{p{\linewidth}}}]{
                \texttt{\tiny{[GM$|$GM]}}
                \texttt{H moved to (3, 4)} \\
            }
        }
    }
    & & \\ \\

    \theutterance \stepcounter{utterance}  
    & & \multicolumn{4}{p{0.6\linewidth}}{
        \cellcolor[rgb]{0.9,0.9,0.9}{
            \makecell[{{p{\linewidth}}}]{
                \texttt{\tiny{[P1$\langle$GM]}}
                \texttt{Veiksmīgi pārvietots 'C' uz (1,4). Jūsu atjauninātais režģis izskatās šādi:} \\
\texttt{```} \\
\texttt{ 1234567} \\
\texttt{╔═╤═╤═══╗} \\
\texttt{║◌│T│◌◌I║ 1} \\
\texttt{║◌├─┤W◌◌║ 2} \\
\texttt{╟─┤◌├─┬─╢ 3} \\
\texttt{║C│◌│◌│◌║ 4} \\
\texttt{╟─┼─┘◌│◌║ 5} \\
\texttt{║◌│◌H┌┴─╢ 6} \\
\texttt{║◌│◌◌│◌◌║ 7} \\
\texttt{╚═╧══╧══╝} \\
\texttt{```} \\
\texttt{Pašlaik jūs spēlējat kārtu 3 no maksimāli 20.} \\
\texttt{Jūs kopā esat uzkrājuši 3 no 12 sodiem.} \\
\texttt{Otrs spēlētājs pārvietoja objektu savā režģī.} \\
\\ 
\\ 
\texttt{Kāda ir jūsu nākamā komanda?} \\
            }
        }
    }
    & \\ \\

    \theutterance \stepcounter{utterance}  
    & \multicolumn{4}{p{0.6\linewidth}}{
        \cellcolor[rgb]{0.9,0.9,0.9}{
            \makecell[{{p{\linewidth}}}]{
                \texttt{\tiny{[P1$\rangle$GM]}}
                \texttt{GĀJIENS: I, (4, 7)} \\
            }
        }
    }
    & & \\ \\

    \theutterance \stepcounter{utterance}  
    & & & \multicolumn{2}{p{0.3\linewidth}}{
        \cellcolor[rgb]{0.9,0.9,0.9}{
            \makecell[{{p{\linewidth}}}]{
                \texttt{\tiny{[GM$|$GM]}}
                \texttt{GĀJIENS: I, (4, 7)} \\
            }
        }
    }
    & & \\ \\

    \theutterance \stepcounter{utterance}  
    & & & \multicolumn{2}{p{0.3\linewidth}}{
        \cellcolor[rgb]{0.9,0.9,0.9}{
            \makecell[{{p{\linewidth}}}]{
                \texttt{\tiny{[GM$|$GM]}}
                \texttt{I moved to (4, 7)} \\
            }
        }
    }
    & & \\ \\

    \theutterance \stepcounter{utterance}  
    & & \multicolumn{4}{p{0.6\linewidth}}{
        \cellcolor[rgb]{0.9,0.9,0.9}{
            \makecell[{{p{\linewidth}}}]{
                \texttt{\tiny{[GM$\rangle$P2]}}
                \texttt{Veiksmīgi pārvietots 'H' uz (3,4). Jūsu atjauninātais režģis izskatās šādi:} \\
\texttt{```} \\
\texttt{ 1234567} \\
\texttt{╔═╤═╤═══╗} \\
\texttt{║◌│◌│T◌◌║ 1} \\
\texttt{║◌├─┤◌◌◌║ 2} \\
\texttt{╟─┤◌├─┬─╢ 3} \\
\texttt{║◌│H│C│◌║ 4} \\
\texttt{╟─┼─┘◌│◌║ 5} \\
\texttt{║◌│◌◌┌┴─╢ 6} \\
\texttt{║◌│WI│◌◌║ 7} \\
\texttt{╚═╧══╧══╝} \\
\texttt{```} \\
\texttt{Pašlaik jūs spēlējat kārtu 4 no maksimāli 20.} \\
\texttt{Jūs kopā esat uzkrājuši 3 no 12 sodiem.} \\
\texttt{Otrs spēlētājs pārvietoja objektu savā režģī.} \\
\\ 
\\ 
\texttt{Kāda ir jūsu nākamā komanda?} \\
            }
        }
    }
    & \\ \\

    \theutterance \stepcounter{utterance}  
    & & & \multicolumn{4}{p{0.6\linewidth}}{
        \cellcolor[rgb]{0.9,0.9,0.9}{
            \makecell[{{p{\linewidth}}}]{
                \texttt{\tiny{[GM$\langle$P2]}}
                \texttt{GĀJIENS: C, (1,4)} \\
            }
        }
    }
     \\ \\

    \theutterance \stepcounter{utterance}  
    & & & \multicolumn{2}{p{0.3\linewidth}}{
        \cellcolor[rgb]{0.9,0.9,0.9}{
            \makecell[{{p{\linewidth}}}]{
                \texttt{\tiny{[GM$|$GM]}}
                \texttt{GĀJIENS: C, (1,4)} \\
            }
        }
    }
    & & \\ \\

    \theutterance \stepcounter{utterance}  
    & & & \multicolumn{2}{p{0.3\linewidth}}{
        \cellcolor[rgb]{0.9,0.9,0.9}{
            \makecell[{{p{\linewidth}}}]{
                \texttt{\tiny{[GM$|$GM]}}
                \texttt{C moved to (1, 4)} \\
            }
        }
    }
    & & \\ \\

    \theutterance \stepcounter{utterance}  
    & & \multicolumn{4}{p{0.6\linewidth}}{
        \cellcolor[rgb]{0.9,0.9,0.9}{
            \makecell[{{p{\linewidth}}}]{
                \texttt{\tiny{[P1$\langle$GM]}}
                \texttt{Veiksmīgi pārvietots 'I' uz (4,7). Jūsu atjauninātais režģis izskatās šādi:} \\
\texttt{```} \\
\texttt{ 1234567} \\
\texttt{╔═╤═╤═══╗} \\
\texttt{║◌│T│◌◌◌║ 1} \\
\texttt{║◌├─┤W◌◌║ 2} \\
\texttt{╟─┤◌├─┬─╢ 3} \\
\texttt{║C│◌│◌│◌║ 4} \\
\texttt{╟─┼─┘◌│◌║ 5} \\
\texttt{║◌│◌H┌┴─╢ 6} \\
\texttt{║◌│◌I│◌◌║ 7} \\
\texttt{╚═╧══╧══╝} \\
\texttt{```} \\
\texttt{Pašlaik jūs spēlējat kārtu 4 no maksimāli 20.} \\
\texttt{Jūs kopā esat uzkrājuši 3 no 12 sodiem.} \\
\texttt{Otrs spēlētājs pārvietoja objektu savā režģī.} \\
\\ 
\\ 
\texttt{Kāda ir jūsu nākamā komanda?} \\
            }
        }
    }
    & \\ \\

    \theutterance \stepcounter{utterance}  
    & \multicolumn{4}{p{0.6\linewidth}}{
        \cellcolor[rgb]{0.9,0.9,0.9}{
            \makecell[{{p{\linewidth}}}]{
                \texttt{\tiny{[P1$\rangle$GM]}}
                \texttt{GĀJIENS: H, (3, 4)} \\
            }
        }
    }
    & & \\ \\

    \theutterance \stepcounter{utterance}  
    & & & \multicolumn{2}{p{0.3\linewidth}}{
        \cellcolor[rgb]{0.9,0.9,0.9}{
            \makecell[{{p{\linewidth}}}]{
                \texttt{\tiny{[GM$|$GM]}}
                \texttt{GĀJIENS: H, (3, 4)} \\
            }
        }
    }
    & & \\ \\

    \theutterance \stepcounter{utterance}  
    & & & \multicolumn{2}{p{0.3\linewidth}}{
        \cellcolor[rgb]{0.9,0.9,0.9}{
            \makecell[{{p{\linewidth}}}]{
                \texttt{\tiny{[GM$|$GM]}}
                \texttt{H moved to (3, 4)} \\
            }
        }
    }
    & & \\ \\

    \theutterance \stepcounter{utterance}  
    & & \multicolumn{4}{p{0.6\linewidth}}{
        \cellcolor[rgb]{0.9,0.9,0.9}{
            \makecell[{{p{\linewidth}}}]{
                \texttt{\tiny{[GM$\rangle$P2]}}
                \texttt{Veiksmīgi pārvietots 'C' uz (1,4). Jūsu atjauninātais režģis izskatās šādi:} \\
\texttt{```} \\
\texttt{ 1234567} \\
\texttt{╔═╤═╤═══╗} \\
\texttt{║◌│◌│T◌◌║ 1} \\
\texttt{║◌├─┤◌◌◌║ 2} \\
\texttt{╟─┤◌├─┬─╢ 3} \\
\texttt{║C│H│◌│◌║ 4} \\
\texttt{╟─┼─┘◌│◌║ 5} \\
\texttt{║◌│◌◌┌┴─╢ 6} \\
\texttt{║◌│WI│◌◌║ 7} \\
\texttt{╚═╧══╧══╝} \\
\texttt{```} \\
\texttt{Pašlaik jūs spēlējat kārtu 5 no maksimāli 20.} \\
\texttt{Jūs kopā esat uzkrājuši 3 no 12 sodiem.} \\
\texttt{Otrs spēlētājs pārvietoja objektu savā režģī.} \\
\\ 
\\ 
\texttt{Kāda ir jūsu nākamā komanda?} \\
            }
        }
    }
    & \\ \\

    \theutterance \stepcounter{utterance}  
    & & & \multicolumn{4}{p{0.6\linewidth}}{
        \cellcolor[rgb]{0.9,0.9,0.9}{
            \makecell[{{p{\linewidth}}}]{
                \texttt{\tiny{[GM$\langle$P2]}}
                \texttt{GĀJIENS: W, (4,6)} \\
            }
        }
    }
     \\ \\

    \theutterance \stepcounter{utterance}  
    & & & \multicolumn{2}{p{0.3\linewidth}}{
        \cellcolor[rgb]{0.9,0.9,0.9}{
            \makecell[{{p{\linewidth}}}]{
                \texttt{\tiny{[GM$|$GM]}}
                \texttt{GĀJIENS: W, (4,6)} \\
            }
        }
    }
    & & \\ \\

    \theutterance \stepcounter{utterance}  
    & & & \multicolumn{2}{p{0.3\linewidth}}{
        \cellcolor[rgb]{0.9,0.9,0.9}{
            \makecell[{{p{\linewidth}}}]{
                \texttt{\tiny{[GM$|$GM]}}
                \texttt{W moved to (4, 6)} \\
            }
        }
    }
    & & \\ \\

    \theutterance \stepcounter{utterance}  
    & & \multicolumn{4}{p{0.6\linewidth}}{
        \cellcolor[rgb]{0.9,0.9,0.9}{
            \makecell[{{p{\linewidth}}}]{
                \texttt{\tiny{[P1$\langle$GM]}}
                \texttt{Veiksmīgi pārvietots 'H' uz (3,4). Jūsu atjauninātais režģis izskatās šādi:} \\
\texttt{```} \\
\texttt{ 1234567} \\
\texttt{╔═╤═╤═══╗} \\
\texttt{║◌│T│◌◌◌║ 1} \\
\texttt{║◌├─┤W◌◌║ 2} \\
\texttt{╟─┤◌├─┬─╢ 3} \\
\texttt{║C│H│◌│◌║ 4} \\
\texttt{╟─┼─┘◌│◌║ 5} \\
\texttt{║◌│◌◌┌┴─╢ 6} \\
\texttt{║◌│◌I│◌◌║ 7} \\
\texttt{╚═╧══╧══╝} \\
\texttt{```} \\
\texttt{Pašlaik jūs spēlējat kārtu 5 no maksimāli 20.} \\
\texttt{Jūs kopā esat uzkrājuši 3 no 12 sodiem.} \\
\texttt{Otrs spēlētājs pārvietoja objektu savā režģī.} \\
\\ 
\\ 
\texttt{Kāda ir jūsu nākamā komanda?} \\
            }
        }
    }
    & \\ \\

    \theutterance \stepcounter{utterance}  
    & \multicolumn{4}{p{0.6\linewidth}}{
        \cellcolor[rgb]{0.9,0.9,0.9}{
            \makecell[{{p{\linewidth}}}]{
                \texttt{\tiny{[P1$\rangle$GM]}}
                \texttt{Looking at my grid, I need to move W from (4,2) to (4,6). Let me check if (4,6) is empty {-} looking at row 6: ║◌│◌◌┌┴─╢, position 4 appears to be empty (◌).} \\
\\ 
\texttt{GĀJIENS: W, (4, 6)} \\
            }
        }
    }
    & & \\ \\

    \theutterance \stepcounter{utterance}  
    & & & \multicolumn{2}{p{0.3\linewidth}}{
        \cellcolor[rgb]{0.9,0.9,0.9}{
            \makecell[{{p{\linewidth}}}]{
                \texttt{\tiny{[GM$|$GM]}}
                \texttt{Looking at my grid, I need to move W from (4,2) to (4,6). Let me check if (4,6) is empty {-} looking at row 6: ║◌│◌◌┌┴─╢, position 4 appears to be empty (◌).} \\
\\ 
\texttt{GĀJIENS: W, (4, 6)} \\
            }
        }
    }
    & & \\ \\

    \theutterance \stepcounter{utterance}  
    & & & \multicolumn{2}{p{0.3\linewidth}}{
        \cellcolor[rgb]{0.9,0.9,0.9}{
            \makecell[{{p{\linewidth}}}]{
                \texttt{\tiny{[GM$|$GM]}}
                \texttt{Invalid format: head is not empty: 'Looking at my grid, I need to move W from (4,2) to (4,6). Let me check if (4,6) is empty {-} looking at row 6: ║◌│◌◌┌┴─╢, position 4 appears to be empty (◌).} \\
\\ 
\texttt{'} \\
            }
        }
    }
    & & \\ \\

    \theutterance \stepcounter{utterance}  
    & & \multicolumn{4}{p{0.6\linewidth}}{
        \cellcolor[rgb]{0.9,0.9,0.9}{
            \makecell[{{p{\linewidth}}}]{
                \texttt{\tiny{[P1$\langle$GM]}}
                \texttt{Sods: Jūsu ziņojumā nedrīkst būt nekas pirms komandas!} \\
\\ 
\texttt{Pārliecinieties, ka jūsu atbildē ir tikai `SAKI: <ZIŅA>` vai `GĀJIENS: <OBJEKTS>, (<X>, <Y>)`, un nekas cits!} \\
\\ 
\texttt{Jūs kopā esat uzkrājuši 4 no 12 sodiem.} \\
\texttt{Lūdzu, mēģiniet vēlreiz!} \\
            }
        }
    }
    & \\ \\

    \theutterance \stepcounter{utterance}  
    & \multicolumn{4}{p{0.6\linewidth}}{
        \cellcolor[rgb]{0.9,0.9,0.9}{
            \makecell[{{p{\linewidth}}}]{
                \texttt{\tiny{[P1$\rangle$GM]}}
                \texttt{GĀJIENS: W, (4, 6)} \\
            }
        }
    }
    & & \\ \\

    \theutterance \stepcounter{utterance}  
    & & & \multicolumn{2}{p{0.3\linewidth}}{
        \cellcolor[rgb]{0.9,0.9,0.9}{
            \makecell[{{p{\linewidth}}}]{
                \texttt{\tiny{[GM$|$GM]}}
                \texttt{GĀJIENS: W, (4, 6)} \\
            }
        }
    }
    & & \\ \\

    \theutterance \stepcounter{utterance}  
    & & & \multicolumn{2}{p{0.3\linewidth}}{
        \cellcolor[rgb]{0.9,0.9,0.9}{
            \makecell[{{p{\linewidth}}}]{
                \texttt{\tiny{[GM$|$GM]}}
                \texttt{W moved to (4, 6)} \\
            }
        }
    }
    & & \\ \\

    \theutterance \stepcounter{utterance}  
    & & \multicolumn{4}{p{0.6\linewidth}}{
        \cellcolor[rgb]{0.9,0.9,0.9}{
            \makecell[{{p{\linewidth}}}]{
                \texttt{\tiny{[GM$\rangle$P2]}}
                \texttt{Veiksmīgi pārvietots 'W' uz (4,6). Jūsu atjauninātais režģis izskatās šādi:} \\
\texttt{```} \\
\texttt{ 1234567} \\
\texttt{╔═╤═╤═══╗} \\
\texttt{║◌│◌│T◌◌║ 1} \\
\texttt{║◌├─┤◌◌◌║ 2} \\
\texttt{╟─┤◌├─┬─╢ 3} \\
\texttt{║C│H│◌│◌║ 4} \\
\texttt{╟─┼─┘◌│◌║ 5} \\
\texttt{║◌│◌W┌┴─╢ 6} \\
\texttt{║◌│◌I│◌◌║ 7} \\
\texttt{╚═╧══╧══╝} \\
\texttt{```} \\
\texttt{Pašlaik jūs spēlējat kārtu 6 no maksimāli 20.} \\
\texttt{Jūs kopā esat uzkrājuši 4 no 12 sodiem.} \\
\texttt{Otrs spēlētājs pārvietoja objektu savā režģī.} \\
\\ 
\\ 
\texttt{Kāda ir jūsu nākamā komanda?} \\
            }
        }
    }
    & \\ \\

    \theutterance \stepcounter{utterance}  
    & & & \multicolumn{4}{p{0.6\linewidth}}{
        \cellcolor[rgb]{0.9,0.9,0.9}{
            \makecell[{{p{\linewidth}}}]{
                \texttt{\tiny{[GM$\langle$P2]}}
                \texttt{GĀJIENS: I, (4,7)} \\
            }
        }
    }
     \\ \\

    \theutterance \stepcounter{utterance}  
    & & & \multicolumn{2}{p{0.3\linewidth}}{
        \cellcolor[rgb]{0.9,0.9,0.9}{
            \makecell[{{p{\linewidth}}}]{
                \texttt{\tiny{[GM$|$GM]}}
                \texttt{GĀJIENS: I, (4,7)} \\
            }
        }
    }
    & & \\ \\

    \theutterance \stepcounter{utterance}  
    & & & \multicolumn{2}{p{0.3\linewidth}}{
        \cellcolor[rgb]{0.9,0.9,0.9}{
            \makecell[{{p{\linewidth}}}]{
                \texttt{\tiny{[GM$|$GM]}}
                \texttt{Sods: (4,7) nav tukšs, bet satur 'I'.} \\
\\ 
\texttt{Jūs kopā esat uzkrājuši 5 no 12 sodiem.} \\
\texttt{Lūdzu, mēģiniet vēlreiz!} \\
            }
        }
    }
    & & \\ \\

    \theutterance \stepcounter{utterance}  
    & & \multicolumn{4}{p{0.6\linewidth}}{
        \cellcolor[rgb]{0.9,0.9,0.9}{
            \makecell[{{p{\linewidth}}}]{
                \texttt{\tiny{[GM$\rangle$P2]}}
                \texttt{Sods: (4,7) nav tukšs, bet satur 'I'.} \\
\\ 
\texttt{Jūs kopā esat uzkrājuši 5 no 12 sodiem.} \\
\texttt{Lūdzu, mēģiniet vēlreiz!} \\
            }
        }
    }
    & \\ \\

    \theutterance \stepcounter{utterance}  
    & & & \multicolumn{4}{p{0.6\linewidth}}{
        \cellcolor[rgb]{0.9,0.9,0.9}{
            \makecell[{{p{\linewidth}}}]{
                \texttt{\tiny{[GM$\langle$P2]}}
                \texttt{Pārbaudīju režģi {-} 'I' jau ir pozīcijā (4,7). Vai vari apstiprināt, ka visi tavi objekti ir vienādās koordinātēs? Manējie ir savās vietās.} \\
\\ 
\texttt{SAKI: pabeigts?} \\
            }
        }
    }
     \\ \\

    \theutterance \stepcounter{utterance}  
    & & & \multicolumn{2}{p{0.3\linewidth}}{
        \cellcolor[rgb]{0.9,0.9,0.9}{
            \makecell[{{p{\linewidth}}}]{
                \texttt{\tiny{[GM$|$GM]}}
                \texttt{Pārbaudīju režģi {-} 'I' jau ir pozīcijā (4,7). Vai vari apstiprināt, ka visi tavi objekti ir vienādās koordinātēs? Manējie ir savās vietās.} \\
\\ 
\texttt{SAKI: pabeigts?} \\
            }
        }
    }
    & & \\ \\

    \theutterance \stepcounter{utterance}  
    & & & \multicolumn{2}{p{0.3\linewidth}}{
        \cellcolor[rgb]{0.9,0.9,0.9}{
            \makecell[{{p{\linewidth}}}]{
                \texttt{\tiny{[GM$|$GM]}}
                \texttt{Invalid format: head is not empty: 'Pārbaudīju režģi {-} 'I' jau ir pozīcijā (4,7). Vai vari apstiprināt, ka visi tavi objekti ir vienādās koordinātēs? Manējie ir savās vietās.} \\
\\ 
\texttt{'} \\
            }
        }
    }
    & & \\ \\

    \theutterance \stepcounter{utterance}  
    & & \multicolumn{4}{p{0.6\linewidth}}{
        \cellcolor[rgb]{0.9,0.9,0.9}{
            \makecell[{{p{\linewidth}}}]{
                \texttt{\tiny{[GM$\rangle$P2]}}
                \texttt{Sods: Jūsu ziņojumā nedrīkst būt nekas pirms komandas!} \\
\\ 
\texttt{Pārliecinieties, ka jūsu atbildē ir tikai `SAKI: <ZIŅA>` vai `GĀJIENS: <OBJEKTS>, (<X>, <Y>)`, un nekas cits!} \\
\\ 
\texttt{Jūs kopā esat uzkrājuši 6 no 12 sodiem.} \\
\texttt{Lūdzu, mēģiniet vēlreiz!} \\
            }
        }
    }
    & \\ \\

    \theutterance \stepcounter{utterance}  
    & & & \multicolumn{4}{p{0.6\linewidth}}{
        \cellcolor[rgb]{0.9,0.9,0.9}{
            \makecell[{{p{\linewidth}}}]{
                \texttt{\tiny{[GM$\langle$P2]}}
                \texttt{SAKI: pabeigts?} \\
            }
        }
    }
     \\ \\

    \theutterance \stepcounter{utterance}  
    & & & \multicolumn{2}{p{0.3\linewidth}}{
        \cellcolor[rgb]{0.9,0.9,0.9}{
            \makecell[{{p{\linewidth}}}]{
                \texttt{\tiny{[GM$|$GM]}}
                \texttt{SAKI: pabeigts?} \\
            }
        }
    }
    & & \\ \\

    \theutterance \stepcounter{utterance}  
    & & \multicolumn{4}{p{0.6\linewidth}}{
        \cellcolor[rgb]{0.9,0.9,0.9}{
            \makecell[{{p{\linewidth}}}]{
                \texttt{\tiny{[P1$\langle$GM]}}
                \texttt{Veiksmīgi pārvietots 'W' uz (4,6). Jūsu atjauninātais režģis izskatās šādi:} \\
\texttt{```} \\
\texttt{ 1234567} \\
\texttt{╔═╤═╤═══╗} \\
\texttt{║◌│T│◌◌◌║ 1} \\
\texttt{║◌├─┤◌◌◌║ 2} \\
\texttt{╟─┤◌├─┬─╢ 3} \\
\texttt{║C│H│◌│◌║ 4} \\
\texttt{╟─┼─┘◌│◌║ 5} \\
\texttt{║◌│◌W┌┴─╢ 6} \\
\texttt{║◌│◌I│◌◌║ 7} \\
\texttt{╚═╧══╧══╝} \\
\texttt{```} \\
\texttt{Pašlaik jūs spēlējat kārtu 6 no maksimāli 20.} \\
\texttt{Jūs kopā esat uzkrājuši 6 no 12 sodiem.} \\
\texttt{Otrs spēlētājs nosūtīja šo ziņojumu:} \\
\\ 
\\ 
\texttt{"pabeigts?"} \\
\\ 
\\ 
\texttt{Kāda ir jūsu nākamā komanda?} \\
            }
        }
    }
    & \\ \\

    \theutterance \stepcounter{utterance}  
    & \multicolumn{4}{p{0.6\linewidth}}{
        \cellcolor[rgb]{0.9,0.9,0.9}{
            \makecell[{{p{\linewidth}}}]{
                \texttt{\tiny{[P1$\rangle$GM]}}
                \texttt{GĀJIENS: T, (5, 1)} \\
            }
        }
    }
    & & \\ \\

    \theutterance \stepcounter{utterance}  
    & & & \multicolumn{2}{p{0.3\linewidth}}{
        \cellcolor[rgb]{0.9,0.9,0.9}{
            \makecell[{{p{\linewidth}}}]{
                \texttt{\tiny{[GM$|$GM]}}
                \texttt{GĀJIENS: T, (5, 1)} \\
            }
        }
    }
    & & \\ \\

    \theutterance \stepcounter{utterance}  
    & & & \multicolumn{2}{p{0.3\linewidth}}{
        \cellcolor[rgb]{0.9,0.9,0.9}{
            \makecell[{{p{\linewidth}}}]{
                \texttt{\tiny{[GM$|$GM]}}
                \texttt{T moved to (5, 1)} \\
            }
        }
    }
    & & \\ \\

    \theutterance \stepcounter{utterance}  
    & & \multicolumn{4}{p{0.6\linewidth}}{
        \cellcolor[rgb]{0.9,0.9,0.9}{
            \makecell[{{p{\linewidth}}}]{
                \texttt{\tiny{[GM$\rangle$P2]}}
                \texttt{Jūsu ziņojums ir nodots otram spēlētājam.} \\
\\ 
\\ 
\texttt{Pašlaik jūs spēlējat kārtu 7 no maksimāli 20.} \\
\texttt{Jūs kopā esat uzkrājuši 6 no 12 sodiem.} \\
\texttt{Otrs spēlētājs pārvietoja objektu savā režģī.} \\
\\ 
\\ 
\texttt{Kāda ir jūsu nākamā komanda?} \\
            }
        }
    }
    & \\ \\

    \theutterance \stepcounter{utterance}  
    & & & \multicolumn{4}{p{0.6\linewidth}}{
        \cellcolor[rgb]{0.9,0.9,0.9}{
            \makecell[{{p{\linewidth}}}]{
                \texttt{\tiny{[GM$\langle$P2]}}
                \texttt{GĀJIENS: I, (4,7)} \\
            }
        }
    }
     \\ \\

    \theutterance \stepcounter{utterance}  
    & & & \multicolumn{2}{p{0.3\linewidth}}{
        \cellcolor[rgb]{0.9,0.9,0.9}{
            \makecell[{{p{\linewidth}}}]{
                \texttt{\tiny{[GM$|$GM]}}
                \texttt{GĀJIENS: I, (4,7)} \\
            }
        }
    }
    & & \\ \\

    \theutterance \stepcounter{utterance}  
    & & & \multicolumn{2}{p{0.3\linewidth}}{
        \cellcolor[rgb]{0.9,0.9,0.9}{
            \makecell[{{p{\linewidth}}}]{
                \texttt{\tiny{[GM$|$GM]}}
                \texttt{Sods: (4,7) nav tukšs, bet satur 'I'.} \\
\\ 
\texttt{Jūs kopā esat uzkrājuši 7 no 12 sodiem.} \\
\texttt{Lūdzu, mēģiniet vēlreiz!} \\
            }
        }
    }
    & & \\ \\

    \theutterance \stepcounter{utterance}  
    & & \multicolumn{4}{p{0.6\linewidth}}{
        \cellcolor[rgb]{0.9,0.9,0.9}{
            \makecell[{{p{\linewidth}}}]{
                \texttt{\tiny{[GM$\rangle$P2]}}
                \texttt{Sods: (4,7) nav tukšs, bet satur 'I'.} \\
\\ 
\texttt{Jūs kopā esat uzkrājuši 7 no 12 sodiem.} \\
\texttt{Lūdzu, mēģiniet vēlreiz!} \\
            }
        }
    }
    & \\ \\

    \theutterance \stepcounter{utterance}  
    & & & \multicolumn{4}{p{0.6\linewidth}}{
        \cellcolor[rgb]{0.9,0.9,0.9}{
            \makecell[{{p{\linewidth}}}]{
                \texttt{\tiny{[GM$\langle$P2]}}
                \texttt{SAKI: Es nevaru pārvietot I, jo tā jau ir (4,7. Vai vari apstiprināt, ka tavi objekti ir mērķa pozīcijās?} \\
            }
        }
    }
     \\ \\

    \theutterance \stepcounter{utterance}  
    & & & \multicolumn{2}{p{0.3\linewidth}}{
        \cellcolor[rgb]{0.9,0.9,0.9}{
            \makecell[{{p{\linewidth}}}]{
                \texttt{\tiny{[GM$|$GM]}}
                \texttt{SAKI: Es nevaru pārvietot I, jo tā jau ir (4,7. Vai vari apstiprināt, ka tavi objekti ir mērķa pozīcijās?} \\
            }
        }
    }
    & & \\ \\

    \theutterance \stepcounter{utterance}  
    & & \multicolumn{4}{p{0.6\linewidth}}{
        \cellcolor[rgb]{0.9,0.9,0.9}{
            \makecell[{{p{\linewidth}}}]{
                \texttt{\tiny{[P1$\langle$GM]}}
                \texttt{Veiksmīgi pārvietots 'T' uz (5,1). Jūsu atjauninātais režģis izskatās šādi:} \\
\texttt{```} \\
\texttt{ 1234567} \\
\texttt{╔═╤═╤═══╗} \\
\texttt{║◌│◌│T◌◌║ 1} \\
\texttt{║◌├─┤◌◌◌║ 2} \\
\texttt{╟─┤◌├─┬─╢ 3} \\
\texttt{║C│H│◌│◌║ 4} \\
\texttt{╟─┼─┘◌│◌║ 5} \\
\texttt{║◌│◌W┌┴─╢ 6} \\
\texttt{║◌│◌I│◌◌║ 7} \\
\texttt{╚═╧══╧══╝} \\
\texttt{```} \\
\texttt{Pašlaik jūs spēlējat kārtu 7 no maksimāli 20.} \\
\texttt{Jūs kopā esat uzkrājuši 7 no 12 sodiem.} \\
\texttt{Otrs spēlētājs nosūtīja šo ziņojumu:} \\
\\ 
\\ 
\texttt{"Es nevaru pārvietot I, jo tā jau ir (4,7. Vai vari apstiprināt, ka tavi objekti ir mērķa pozīcijās?"} \\
\\ 
\\ 
\texttt{Kāda ir jūsu nākamā komanda?} \\
            }
        }
    }
    & \\ \\

    \theutterance \stepcounter{utterance}  
    & \multicolumn{4}{p{0.6\linewidth}}{
        \cellcolor[rgb]{0.9,0.9,0.9}{
            \makecell[{{p{\linewidth}}}]{
                \texttt{\tiny{[P1$\rangle$GM]}}
                \texttt{SAKI: pabeigts!} \\
            }
        }
    }
    & & \\ \\

    \theutterance \stepcounter{utterance}  
    & & & \multicolumn{2}{p{0.3\linewidth}}{
        \cellcolor[rgb]{0.9,0.9,0.9}{
            \makecell[{{p{\linewidth}}}]{
                \texttt{\tiny{[GM$|$GM]}}
                \texttt{SAKI: pabeigts!} \\
            }
        }
    }
    & & \\ \\

    \theutterance \stepcounter{utterance}  
    & & & \multicolumn{2}{p{0.3\linewidth}}{
        \cellcolor[rgb]{0.9,0.9,0.9}{
            \makecell[{{p{\linewidth}}}]{
                \texttt{\tiny{[GM$|$GM]}}
                \texttt{true} \\
            }
        }
    }
    & & \\ \\

    \theutterance \stepcounter{utterance}  
    & & & \multicolumn{2}{p{0.3\linewidth}}{
        \cellcolor[rgb]{0.9,0.9,0.9}{
            \makecell[{{p{\linewidth}}}]{
                \texttt{\tiny{[GM$|$GM]}}
                \texttt{Initial states:} \\
\texttt{```} \\
\texttt{Player 1:} \\
\texttt{ 1234567} \\
\texttt{╔═╤═╤═══╗} \\
\texttt{║C│◌│◌TI║ 1} \\
\texttt{║◌├─┤W◌◌║ 2} \\
\texttt{╟─┤◌├─┬─╢ 3} \\
\texttt{║◌│◌│◌│◌║ 4} \\
\texttt{╟─┼─┘◌│◌║ 5} \\
\texttt{║◌│◌H┌┴─╢ 6} \\
\texttt{║◌│◌◌│◌◌║ 7} \\
\texttt{╚═╧══╧══╝} \\
\\ 
\\ 
\texttt{Player 2:} \\
\texttt{ 1234567} \\
\texttt{╔═╤═╤═══╗} \\
\texttt{║H│◌│◌◌◌║ 1} \\
\texttt{║◌├─┤◌◌◌║ 2} \\
\texttt{╟─┤◌├─┬─╢ 3} \\
\texttt{║◌│◌│C│◌║ 4} \\
\texttt{╟─┼─┘◌│◌║ 5} \\
\texttt{║◌│◌◌┌┴─╢ 6} \\
\texttt{║T│WI│◌◌║ 7} \\
\texttt{╚═╧══╧══╝} \\
\\ 
\texttt{```} \\
            }
        }
    }
    & & \\ \\

    \theutterance \stepcounter{utterance}  
    & & & \multicolumn{2}{p{0.3\linewidth}}{
        \cellcolor[rgb]{0.9,0.9,0.9}{
            \makecell[{{p{\linewidth}}}]{
                \texttt{\tiny{[GM$|$GM]}}
                \texttt{End states:} \\
\texttt{```} \\
\texttt{Player 1:} \\
\texttt{ 1234567} \\
\texttt{╔═╤═╤═══╗} \\
\texttt{║◌│◌│T◌◌║ 1} \\
\texttt{║◌├─┤◌◌◌║ 2} \\
\texttt{╟─┤◌├─┬─╢ 3} \\
\texttt{║C│H│◌│◌║ 4} \\
\texttt{╟─┼─┘◌│◌║ 5} \\
\texttt{║◌│◌W┌┴─╢ 6} \\
\texttt{║◌│◌I│◌◌║ 7} \\
\texttt{╚═╧══╧══╝} \\
\\ 
\\ 
\texttt{Player 2:} \\
\texttt{ 1234567} \\
\texttt{╔═╤═╤═══╗} \\
\texttt{║◌│◌│T◌◌║ 1} \\
\texttt{║◌├─┤◌◌◌║ 2} \\
\texttt{╟─┤◌├─┬─╢ 3} \\
\texttt{║C│H│◌│◌║ 4} \\
\texttt{╟─┼─┘◌│◌║ 5} \\
\texttt{║◌│◌W┌┴─╢ 6} \\
\texttt{║◌│◌I│◌◌║ 7} \\
\texttt{╚═╧══╧══╝} \\
\\ 
\texttt{```} \\
            }
        }
    }
    & & \\ \\

    \theutterance \stepcounter{utterance}  
    & & & \multicolumn{2}{p{0.3\linewidth}}{
        \cellcolor[rgb]{0.9,0.9,0.9}{
            \makecell[{{p{\linewidth}}}]{
                \texttt{\tiny{[GM$|$GM]}}
                \texttt{* success: True} \\
\texttt{* lose: False} \\
\texttt{* aborted: False} \\
\texttt{{-}{-}{-}{-}{-}{-}{-}} \\
\texttt{* Shifts: 7.00} \\
\texttt{* Max Shifts: 8.00} \\
\texttt{* Min Shifts: 4.00} \\
\texttt{* End Distance Sum: 0.00} \\
\texttt{* Init Distance Sum: 30.73} \\
\texttt{* Expected Distance Sum: 20.95} \\
\texttt{* Penalties: 7.00} \\
\texttt{* Max Penalties: 12.00} \\
\texttt{* Rounds: 7.00} \\
\texttt{* Max Rounds: 20.00} \\
\texttt{* Object Count: 5.00} \\
            }
        }
    }
    & & \\ \\

\end{supertabular}
}

\end{document}
